%% ---------------------------------------------------------------------------
\section{Migration order and validation}
\label{sec:findings}

Cost is one half of a migration plan; the other is what the migrated credential
protects. This section asks which partial migrations leave which property of
Table~\ref{tab:properties} forgeable, and whether the reference implementation enforces the
rules a migration relies on.

\subsection{When the chain does not share one algorithm}
\label{sec:mixedchain}

Every chain so far uses one algorithm from leaf to root, and no migration
looks like that. A certificate carries two algorithms: the key it holds and
the signature its issuer put on it. Those come apart during a transition, and
in a particular order, because leaf certificates are reissued in days,
intermediates live for years, and roots are pinned in trust stores and move
last.

\begin{table}[t]
\caption{Chains whose levels migrate at different times, leaf first and root last; the root signs the intermediate but is not carried. \emph{Signer PQ} is whether every load-bearing signature in the signer's path is post-quantum. The last two columns add an RFC~3161 timestamp: with a classical TSA the time evidence stays forgeable, so a quantum adversary can backdate the manifest even when the signer's path is post-quantum.}
\label{tab:mixedchain}
\small
\setlength{\tabcolsep}{3.5pt}
\begin{tabular}{llrcrr}
\toprule
 & Chain & COSE & Signer & \multicolumn{2}{c}{+ timestamp (B)}\\
\cmidrule(lr){5-6}
Stage & leaf/int./root & (B) & PQ & ES256 TSA & ML-DSA-65 TSA\\
\midrule
Classical & ES256/ES256/ES256 & 1,585 & no & 3,651 & 17,209\\
Leaf migrated & D65/ES256/ES256 & 6,724 & no & 8,788 & 22,348\\
Leaf and intermediate & D65/D65/ES256 & 11,867 & no & 13,933 & 27,491\\
PKI first, leaves last & ES256/D65/D65 & 9,971 & no & -- & --\\
Fully migrated & D65/D65/D65 & 15,113 & yes & 17,179 & 30,737\\
Assurance ladder & D44/D65/D87 & 14,902 & yes & -- & --\\
Ladder inverted & D87/D65/D44 & 16,182 & yes & -- & --\\
Cheap leaves, strong root & D44/D44/D87 & 13,373 & yes & -- & --\\
\bottomrule
\end{tabular}
\end{table}


Table~\ref{tab:mixedchain} separates two things usually spoken of as one.
Bytes are set by which levels have migrated. Quantum resistance is set by the
weakest load-bearing signature: the manifest signature, the signature on the
leaf certificate and the signature on the intermediate. The root's
self-signature is not load-bearing, and the root is not carried, but its
algorithm still matters because it signs the intermediate. A post-quantum
leaf under the existing classical PKI costs \NmixLeafMig{} times the
classical structure, \NmixPctFullLeafMig\% of what full migration costs;
migrating the intermediate too brings it to \NmixLeafInt{} times. Neither
makes the signer's path quantum resistant, because an adversary who can forge
the issuer's ECDSA signature mints a leaf certificate and signs whatever it
likes.

Migrating the PKI first while issuing classical leaves for compatibility is
the order an operator reaches by reasoning about disruption rather than about
adversaries. It costs \NmixPkiFirst{} times the classical structure,
\NmixPctFullPkiFirst\% of the full price, and leaves the manifest itself
forgeable, since its signature is still ECDSA. What it does buy is a
post-quantum trust infrastructure ready to issue post-quantum leaves without a
second change at the anchor, so its value lies in what follows it, not in any
protection it gives the manifests signed meanwhile.

The timestamp needs its own migration. With a classical time-stamping
authority, a fully migrated signer's manifest can still be backdated by an
adversary who forges the timestamp, which undoes the argument a timestamp
exists to make, that the manifest was signed before its certificate was
revoked. A post-quantum timestamp costs \NmixWholeFullPq{} bytes of manifest
cryptography against \NmixWholeFullEs{} with an ES256 authority, and the
provenance is quantum resistant only when both are post-quantum.

An assurance ladder, a category-2 leaf under a category-3 intermediate under a
category-5 root, puts the higher margin on the longer-lived key and costs
\NmixLadder{} times the classical structure, against \NmixFull{} for a uniform
category-3 chain: the strongest key sits in the root, which is not carried.
Inverting the ladder costs \NmixInverted{} times and is the configuration to
check for in review, since it puts the weakest parameter set on the key that
lives longest and is hardest to replace.


\subsection{What the reference implementation accepts}
\label{sec:validation}

\begin{table*}[t]
\caption{Validation through the reference implementation (\texttt{c2pa-rs} \texttt{e2ad88e} with our ML-DSA patches, p2): a selection of the 40 cases; Table~\ref{tab:valfull} in the supplement lists all. T, V, I: Trusted, Valid, Invalid; E: reader error; SR: signing refused (\emph{direct}: our harness's own structure with that certificate). \emph{Expected} follows C2PA~2.5. \emph{TST}: the timestamp's status (t trusted, m mismatch, u untrusted); an ignored timestamp leaves the signer judged now. The TSA is ES256 except in the indented rows. $^\dagger$Bold: a defect, also present unpatched. The CA-signed \emph{revoked} OCSP row is where C2PA~2.5 and RFC~6960 disagree (Section~\ref{sec:validation}).}
\label{tab:validation}
\footnotesize
\setlength{\tabcolsep}{3pt}
\begin{tabular}{llccc}
\toprule
 & & ES256, & ML-DSA & \\
Case & Expected & Ed25519 & 44/65/87 & TST\\
\midrule
Leaf and intermediate & T & T & T & --\\
\quad timestamped & T & T & T & t\\
Leaf, intermediate, root & T & T & T & t\\
Leaf only, issuer is anchor & T & T & T & t\\
Leaf only, intermediate missing & not T & V & V & t\\
Unknown anchor & not T & V & V & t\\
Thumbprint (\texttt{x5t}) & not T & E & E & --\\
Leaf EKU serverAuth & refused & SR (direct I) & SR (direct I) & -- (direct --)\\
Leaf EKU timeStamping only$^\dagger$ & not T & \textbf{T} & \textbf{T} & t\\
COSE algorithm $\neq$ key (ML-DSA-44 only) & I & I & I & --\\
Image changed after signing & I & I & I & t\\
Signature bit flipped & I & I & I & m\\
Expired leaf, timestamp inside validity & T & T & T & t\\
Expired leaf, no timestamp & not T & I & I & --\\
Expired leaf, timestamp over other data & not T & I & I & m\\
Expired leaf, TSA not trusted & not T & I & I & u\\
Timestamp over other data & T & T & T & m\\
TSA not trusted & T & T & T & u\\
TSA certificate lacks timeStamping$^\dagger$ & T & \textbf{I} & \textbf{I} & u\\
ML-DSA-65 TSA & T & T & T & t\\
\quad timestamp over other data & T & T & T & m\\
\quad expired leaf, timestamp inside validity & T & T & T & t\\
OCSP, CA-signed, revoked & T & T & T & t\\
OCSP, delegated responder, revoked & not T & I & I & t\\
\bottomrule
\end{tabular}
\end{table*}


Table~\ref{tab:validation} gives a selection of the \Nvalcases{} cases of
Section~\ref{sec:patch}; Table~\ref{tab:valfull} in the supplement lists every
case for every algorithm. With both patches every ML-DSA case behaves as the
classical ones do, and with the signing patch alone none reaches
\emph{Trusted}, because path validation cannot check an ML-DSA certificate.

\paragraph{A fully post-quantum credential works}
A credential with an ES256 time-stamping authority validates in
\NvalVerifyLo{} to \NvalVerifyHi\,ms on the M1 Pro. An ML-DSA-65 signer,
timestamped by an ML-DSA-65 time-stamping authority under an ML-DSA-65 anchor,
validates as \emph{Trusted} with a trusted timestamp in \NvalPqTsaMs\,ms, and still does after the signer's certificate
has expired, which is the job a timestamp exists to do. A timestamp over other
data, or from an untrusted authority, is ignored as the specification
says~\cite{c2pa25}, and an expired signer with such a timestamp is rejected, so
a rejected timestamp never lends a signer historical validity. The library
rejects a modified image, a flipped signature bit and a COSE algorithm that
does not match the key, stops short of \emph{Trusted} for a chain to an unknown
anchor or a leaf whose intermediate is missing, refuses to sign with a leaf
that fails its own certificate profile, and accepts the leaf alone when the
leaf's issuer is an anchor, the configuration Section~\ref{sec:carrying}
recommends. It does the same for the time-stamping authority: a token that
carries only the authority's certificate yields a trusted timestamp when the
authority's issuing CA is a time-stamping anchor, for every algorithm, and an
untrusted one when only the root is. A manifest carrying a thumbprint fails with a reader error rather
than a validation status, a 50-certificate chain is reported \emph{Valid} but not \emph{Trusted}, and a
1\,MB certificate validates in \NvalBigCertMs\,ms using \NvalBigCertMB\,MB of
memory.

\paragraph{A time-stamping certificate is accepted as a claim signer}
A leaf whose only extended key usage is \texttt{id-kp-timeStamping} validates
as a \emph{Trusted} claim signer, for every algorithm, and still does under a
policy whose usage list names only \texttt{c2pa-kp-claimSigning}. C2PA~2.5
says that a certificate valid for time-stamping shall be valid for that purpose
alone~\cite{c2pa25}. The library's key-usage check accepts e-mail protection,
time-stamping and OCSP signing unconditionally, whatever the configured list
says. Two added cases confirm it for every algorithm, through the tagged
release for ES256 and Ed25519 and through the patched library for ML-DSA: a
leaf valid only for e-mail protection, or only for OCSP signing, also
validates as a \emph{Trusted} claim signer under the claim-signing-only
policy, and the OCSP-signing leaf does so even though C2PA~2.5 confines such a
certificate to that one purpose. Whoever holds the key of a time-stamping
authority whose certificate chains to a claim-signing anchor can therefore
sign manifests that a validator using the library's key-usage check reports as
\emph{Trusted}. Such a manifest is attributed to the authority's own subject
name, not to an arbitrary publisher: the defect lets a key operated for
another purpose, often by another party, produce trusted provenance, which
the separation the specification draws is meant to prevent, but it does not by
itself let that key impersonate another signer. On the live lists the
condition is not hypothetical, and we checked it with real certificates. We
pulled the time-stamping certificates out of the RFC 3161 tokens of the
production files (Section~\ref{sec:prodchains}), \NtpcChains{} distinct chains
in all, and ran each through the library's own certificate-profile and trust
checks against the C2PA claim-signing list, with no signature. Both checks pass
for \NtpcStrong{} of the \NtpcChains{} chains as a claim signer under a
claim-signing-only policy: \NtpcGoogle{} of Google's time-stamping leaves and
vivo's, the ones whose issuing CA is on the C2PA time-stamping list and is
itself issued by a claim-signing anchor. These \NtpcStrong{} chains sign the
timestamps in \NtpcFiles{} of the production files. For each, the profile
check accepts the time-stamping-only EKU and the trust check builds a path to a
claim-signing anchor, so the library would report a manifest signed by the
key behind any of them as \emph{Trusted}; only the private key is absent. The
remaining chains fail because their issuer is not on the claim-signing list
(DigiCert's, Sectigo's and other general-purpose time-stamping authorities) or,
for \NtpcCaRej{} of them, because the profile check correctly refuses a CA
certificate. The two trust lists make the overlap structural rather than
incidental: \NtlBothLists{} of the \NtlClaimRoots{} roots on the claim-signing
list are also roots on the time-stamping list, and all \NtsaIcaUnderClaim{}
time-stamping issuing CAs on the time-stamping list are issued by claim-signing
roots, so any time-stamping authority listed under one of them is in the same
position as Google's. The behaviour reproduces in the unmodified library at commit
\texttt{\NvalCommit}, in the latest tagged release, c2pa-v0.91.9, and in the
command-line tool c2patool~0.28.4, so it is not an artefact of our patch;
Table~\ref{tab:evidence} in the supplement records the production-certificate
check and what else was reproduced where.

\paragraph{A timestamp the validator should ignore invalidates the manifest}
In the other direction, a time-stamping authority whose certificate lacks the
time-stamping usage makes the whole manifest \emph{Invalid}, with status
\texttt{signingCredential.invalid}, where the specification says to report the
timestamp as untrusted and judge the signer at the current time. The token
travels in the COSE unprotected header, so anyone who can rewrite a file can
swap in such a token, and a misconfigured authority produces one by accident.
Either way valid content is reported as \emph{Invalid}, the status of
tampered content, instead of falling back to validation at the current time.

\paragraph{Revocation stapled in a CA-signed response has no effect}
Unlike the two defects above, this behaviour arises where the C2PA profile and
RFC~6960 disagree and the implementation follows the profile.
RFC~6960 lets the issuing CA sign an OCSP response itself, with no special key
usage~\cite{rfc6960}. C2PA~2.5 instead requires a response to be signed by an
end-entity certificate valid for \texttt{id-kp-OCSPSigning}, while also asking
validators to check stapled responses as RFC~6960 does, which accepts the
CA-signed form~\cite{c2pa25}. The reference implementation resolves the
tension by ignoring the response: it looks for the responder's certificate in
the response, finds none, and treats the status as unknown, the same outcome as
when nothing is stapled. A signer whose stapled, CA-signed response says
\emph{revoked} therefore validates as \emph{Trusted}
(Table~\ref{tab:validation}), with online OCSP fetching off as it was for the
whole matrix; a validator that also fetches status could still catch it. The risk is to
deployments that believe their stapled evidence protects them. A generator that
staples revocation status, as \NcorpProdOcsp{} production signatures do, must
use a delegated responder for it to count; under ML-DSA-65 a delegated
response costs \NocspDelDSAsixtyfive{} bytes per manifest, because it carries
the responder's certificate, against \NocspCADSAsixtyfive{} for a CA-signed
one. Golaszewski et al.\ found that
revocation checking is optional in C2PA~\cite{golaszewski2026}; here evidence
is present and ignored.

\paragraph{Disclosure}
I reported the invalidating-timestamp behaviour publicly, with minimal
vectors, as c2pa-rs issue \#3005 on 3 November 2026~\cite{c2pars3005}. The
key-usage behaviour lets a time-stamping key sign trusted manifests, so I
reported it privately, with its vectors, to Adobe's product security team,
which handles security reports for the library, on 9 November 2026. Its cause
and its reach are described here, but its signed vectors are withheld from the
released data until a fix ships or 90 days after the report, whichever comes
first; its status codes are released. The OCSP behaviour follows the
specification, so I raised it with the specification's
editors~\cite{c2paspec143}, as I did the reading behind leaf-only
carriage~\cite{c2paspec142}.

\paragraph{Reading the findings}
The findings are two implementation defects, the key-usage and
invalidating-timestamp behaviours, and one interoperability and
revocation-policy issue, the CA-signed OCSP response that the C2PA profile and
RFC~6960 treat differently. None is specific to post-quantum algorithms; each
appears with ES256 and Ed25519. Each, however, breaks a property a migration
relies on: that a time-stamping key cannot attribute content, that a timestamp
is judged as the specification says, and that a revoked signer stays revoked.
Several of the Content Authenticity Initiative's SDKs for other languages reuse
this library, so the findings may affect those integrations too, though we did
not evaluate any application built on it or measure deployment share. The two
defects are of the kind that adversarial testing of TLS libraries exposed a
decade ago~\cite{brubaker2014frankencerts,debnath2021ceres}. They surfaced
because adding an algorithm called for a conformance matrix with negative
cases, and a migration that adds the algorithm without one carries such defects
over unnoticed.
